\documentclass[11pt]{article}

\usepackage[margin=1in]{geometry}
\usepackage{amsmath,amssymb,amsthm,mathtools}
\usepackage[T1]{fontenc}
\usepackage{lmodern}
\usepackage{microtype}
\usepackage{hyperref}

\numberwithin{equation}{section}

\newtheorem{theorem}{Theorem}
\newtheorem{proposition}{Proposition}
\newtheorem{remark}{Remark}
\newcommand{\E}{\mathcal E}
\newcommand{\Om}{\Omega}
\newcommand{\dd}{\,\mathrm d}
\newcommand{\Ltwo}{L^2(\Om)}
\newcommand{\norm}[1]{\left\lVert #1\right\rVert}

\title{E14: The Band-Pass Lyapunov Functional Generates the E13 Flow\\
\large A Complete Variational and Energy-Dissipation Proof}
\date{}

\begin{document}
\maketitle

\begin{abstract}
We prove that the band-pass energy
\[
\E[A]=\int_\Om\left(
-r|A|^2-a|\nabla A|^2+b|\Delta A|^2
+\frac{\beta}{2}|A|^4
\right)\dd x
\]
generates the nonlinear band-pass amplitude equation
\[
\partial_tA=(r-a\Delta-b\Delta^2)A-\beta|A|^2A
\]
as its negative complex $L^2$ gradient flow. Along every sufficiently
regular solution, the exact identity
\[
\frac{\dd}{\dd t}\E[A(t)]=-2\norm{\partial_t A(t)}_{L^2}^2
\]
holds. We establish the exact equivalence with the centered
Swift--Hohenberg energy, identify the linear band-pass symbol, and
state the conditions for lower boundedness. The theorem concerns the
declared dissipative amplitude model; it does not derive that model
from a separate higher-time physical action.
\end{abstract}

\section{Setting and conventions}

Let $\Om=\mathbb T^d$ be a flat $d$-dimensional torus with periodic
boundary conditions, where $1\le d\le 3$. Suppose
\[
r\in\mathbb R,\qquad a,b>0,\qquad \beta\ge0.
\]
Let $A:\Om\to\mathbb C$. The energy functional is
\begin{equation}\label{eq:E14}
\boxed{\displaystyle
\E[A]=\int_\Om\left[
-r|A|^2-a|\nabla A|^2+b|\Delta A|^2
+\frac{\beta}{2}|A|^4
\right]\dd x.}
\end{equation}
The expression is well defined on $H^2(\Om;\mathbb C)$, since in
dimensions $d\le3$ the embedding
$H^2(\Om)\hookrightarrow L^\infty(\Om)$ holds.

We define the variational derivative with respect to $\overline A$ by
the convention
\begin{equation}\label{eq:convention}
D\E[A](h)=2\operatorname{Re}\int_\Om
\frac{\delta\E}{\delta\overline A}\,\overline h\dd x
\end{equation}
for smooth complex variations $h$. Thus the negative complex
$L^2$ gradient flow is
\[
\partial_t A=-\frac{\delta\E}{\delta\overline A}.
\]
With the alternative \emph{real} Hilbert-space convention
$(u,v)_{\mathbb R}=\operatorname{Re}\int_\Om
u\overline v\dd x$, the real gradient is
$2\,\delta\E/\delta\overline A$; the same evolution is the
negative real gradient flow of $\E/2$.

\begin{theorem}[E14 gradient-flow identity]\label{thm:main}
For every smooth $A$ and every smooth variation $h$,
\begin{equation}\label{eq:var}
\boxed{\displaystyle
\frac{\delta\E}{\delta\overline A}
=-rA+a\Delta A+b\Delta^2A+\beta|A|^2A.}
\end{equation}
Consequently the negative complex $L^2$ gradient flow of
$\E$ is exactly
\begin{equation}\label{eq:E13}
\boxed{\displaystyle
\partial_t A=(r-a\Delta-b\Delta^2)A-\beta|A|^2A.}
\end{equation}
If $A\in C^1([0,T];H^4(\Om;\mathbb C))$ satisfies
\eqref{eq:E13}, then
\begin{equation}\label{eq:diss}
\boxed{\displaystyle
\frac{\dd}{\dd t}\E[A(t)]
=-2\norm{\partial_t A(t)}_{L^2}^2\le0.}
\end{equation}
In particular, for $0\le t_1\le t_2\le T$,
\begin{equation}\label{eq:integrated}
\boxed{\displaystyle
\E[A(t_2)]+2\int_{t_1}^{t_2}
\norm{\partial_t A(t)}_{L^2}^2\dd t
=\E[A(t_1)].}
\end{equation}
At any time where the solution is differentiable,
the energy derivative vanishes if and only if
$\partial_t A=0$ in $L^2(\Om)$.
\end{theorem}

\begin{proof}
Fix a smooth complex perturbation $h$, let
$A_\varepsilon=A+\varepsilon h$, and differentiate
with respect to the \emph{real} parameter $\varepsilon$
at $\varepsilon=0$.

First,
\begin{equation}\label{eq:var1}
\left.\frac{\dd}{\dd\varepsilon}\right|_0
\bigl(-r|A_\varepsilon|^2\bigr)
=-2r\operatorname{Re}(A\overline h).
\end{equation}
Second,
\begin{equation}\label{eq:var2}
\left.\frac{\dd}{\dd\varepsilon}\right|_0
\bigl(-a|\nabla A_\varepsilon|^2\bigr)
=-2a\operatorname{Re}
(\nabla A\cdot\nabla\overline h).
\end{equation}
Third,
\begin{equation}\label{eq:var3}
\left.\frac{\dd}{\dd\varepsilon}\right|_0
\bigl(b|\Delta A_\varepsilon|^2\bigr)
=2b\operatorname{Re}(\Delta A\,\Delta\overline h).
\end{equation}
Finally, since
\[
\left.\frac{\dd}{\dd\varepsilon}\right|_0
|A+\varepsilon h|^4
=4|A|^2\operatorname{Re}(A\overline h),
\]
we have
\begin{equation}\label{eq:var4}
\left.\frac{\dd}{\dd\varepsilon}\right|_0
\frac{\beta}{2}|A_\varepsilon|^4
=2\beta|A|^2\operatorname{Re}(A\overline h).
\end{equation}
Summing these variations,
\begin{align}
D\E[A](h)
=2\operatorname{Re}\int_\Om\bigl[
&-rA\overline h
-a\nabla A\cdot\nabla\overline h \notag\\
&+b\Delta A\,\Delta\overline h
+\beta|A|^2A\overline h\bigr]\dd x.
\label{eq:firstvariation}
\end{align}

Periodicity eliminates all boundary terms. In particular,
\begin{align*}
\int_\Om\nabla A\cdot\nabla\overline h\dd x
&=-\int_\Om\Delta A\,\overline h\dd x,\\
\int_\Om\Delta A\,\Delta\overline h\dd x
&=\int_\Om\Delta^2 A\,\overline h\dd x.
\end{align*}
It follows that
\begin{equation}\label{eq:variationfinal}
D\E[A](h)=2\operatorname{Re}\int_\Om
\left[-rA+a\Delta A+b\Delta^2 A
+\beta|A|^2 A\right]\overline h\dd x.
\end{equation}
By \eqref{eq:convention} and arbitrariness of $h$,
this proves \eqref{eq:var}. Its negative is
\[
-\frac{\delta\E}{\delta\overline A}
=rA-a\Delta A-b\Delta^2A-\beta|A|^2A,
\]
which proves \eqref{eq:E13}.

Now let $A(t)$ be a solution with the stated regularity.
The chain rule and \eqref{eq:variationfinal} yield
\begin{align*}
\frac{\dd}{\dd t}\E[A(t)]
&=D\E[A(t)](\partial_t A)\\
&=2\operatorname{Re}\int_\Om
\frac{\delta\E}{\delta\overline A}
\,\overline{\partial_t A}\dd x\\
&=-2\operatorname{Re}\int_\Om
\partial_t A\,\overline{\partial_t A}\dd x\\
&=-2\norm{\partial_t A}_{L^2}^2.
\end{align*}
This establishes \eqref{eq:diss}.
Integration in time gives \eqref{eq:integrated}.
Because the integrand in the squared norm is nonnegative,
equality in \eqref{eq:diss} occurs exactly when
$\partial_t A=0$ almost everywhere in space.
\end{proof}

\section{Centered Swift--Hohenberg representation}

\begin{proposition}[Exact complete-square equivalence]
Let
\begin{equation}\label{eq:centerparams}
k_\star^2=\frac{a}{2b},
\qquad
\mu=r+\frac{a^2}{4b},
\qquad
L=\Delta+k_\star^2.
\end{equation}
Then \eqref{eq:E14} equals the centered energy
\begin{equation}\label{eq:centerE}
\boxed{\displaystyle
\E[A]=\int_\Om\left[
-\mu|A|^2+b|LA|^2+\frac{\beta}{2}|A|^4
\right]\dd x.}
\end{equation}
Its negative complex gradient flow is equivalently
\begin{equation}\label{eq:centerflow}
\boxed{\displaystyle
\partial_tA=\mu A-b(\Delta+k_\star^2)^2A
-\beta|A|^2A.}
\end{equation}
\end{proposition}

\begin{proof}
Expanding the squared differential expression,
\[
|LA|^2=|\Delta A|^2+
2k_\star^2\operatorname{Re}(\Delta A\,\overline A)
+k_\star^4|A|^2.
\]
Periodic integration by parts gives
\[
\operatorname{Re}\int_\Om
\Delta A\,\overline A\dd x
=-\int_\Om|\nabla A|^2\dd x.
\]
Consequently, the right-hand side of \eqref{eq:centerE}
is
\begin{align*}
\int_\Om\bigl[
&b|\Delta A|^2-2bk_\star^2|\nabla A|^2\\
&+(bk_\star^4-\mu)|A|^2
+\tfrac{\beta}{2}|A|^4\bigr]\dd x.
\end{align*}
By \eqref{eq:centerparams},
\[
2bk_\star^2=a,
\qquad
bk_\star^4-\mu
=\frac{a^2}{4b}
-\left(r+\frac{a^2}{4b}\right)=-r.
\]
This is exactly \eqref{eq:E14}. Since $L$ is self-adjoint,
\[
\frac{\delta\E}{\delta\overline A}
=-\mu A+bL^2A+\beta|A|^2A.
\]
Taking the negative yields \eqref{eq:centerflow}.
Expanding $L^2$ and using \eqref{eq:centerparams}
recovers \eqref{eq:E13}.
\end{proof}

\section{Linear spectral band selection}

Linearizing \eqref{eq:E13} at $A=0$ gives
\[
\partial_tA=(r-a\Delta-b\Delta^2)A.
\]
For a Fourier mode with spatial frequency $k$,
$\widehat{\Delta A}(k)=-|k|^2\widehat A(k)$, so
\begin{equation}\label{eq:lambda}
\partial_t\widehat A(k)
=\lambda(k)\widehat A(k),\qquad
\lambda(k)=r+a|k|^2-b|k|^4.
\end{equation}
Completing the square,
\begin{equation}\label{eq:bandshell}
\boxed{\displaystyle
\lambda(k)=\mu-b(|k|^2-k_\star^2)^2.}
\end{equation}
Since $b>0$, the continuous symbol is maximized on the
shell $|k|=k_\star$. On a finite torus, the permitted
frequencies are discrete, so an exact mode on that
shell need not exist. As $|k|\to\infty$,
$\lambda(k)\to-\infty$, establishing ultraviolet linear
damping for the \emph{negative} fourth-order sign.

\section{Lower bound and finite total dissipation}

\begin{proposition}[Energy lower bound]
Suppose $\beta>0$, and put $\mu_+=\max\{\mu,0\}$.
For all $A\in H^2(\Om;\mathbb C)$,
\begin{equation}\label{eq:lower}
\boxed{\displaystyle
\E[A]\ge
-\frac{\mu_+^2}{2\beta}|\Om|.}
\end{equation}
If a sufficiently regular solution exists for all $t\ge0$,
then
\begin{equation}\label{eq:infinite}
\boxed{\displaystyle
\int_0^\infty\norm{\partial_tA(t)}_{L^2}^2\dd t
\le\frac12\left[
\E[A(0)]
+\frac{\mu_+^2}{2\beta}|\Om|
\right].}
\end{equation}
\end{proposition}

\begin{proof}
The centered representation gives
\[
\E[A]=b\norm{LA}_{L^2}^2+
\int_\Om\left[
\frac{\beta}{2}|A|^4-\mu|A|^2
\right]\dd x.
\]
For $s\ge0$, we have
\[
\frac{\beta}{2}s^2-\mu s
\ge\frac{\beta}{2}s^2-\mu_+s
=\frac{\beta}{2}
\left(s-\frac{\mu_+}{\beta}\right)^2
-\frac{\mu_+^2}{2\beta}
\ge-\frac{\mu_+^2}{2\beta}.
\]
Using $s=|A|^2$ and $b\norm{LA}_{L^2}^2\ge0$
proves \eqref{eq:lower}. Combining it with
\eqref{eq:integrated}, taking $t_1=0$, and then
sending $t_2\to\infty$ proves \eqref{eq:infinite}.
\end{proof}

\begin{remark}[Meaning of strict dissipation]
Any periodic-in-time classical solution must be stationary:
over a complete temporal period, the left-hand side of
\eqref{eq:integrated} has equal endpoint energies, forcing
$\partial_t A=0$ throughout the period. Nevertheless,
finite total dissipation does \emph{not} by itself prove
convergence to a unique equilibrium.
\end{remark}

\section{Scope and limitations}

\begin{enumerate}
\item The first-variation and dissipation identities are
exact for the declared E13 evolution. Their algebraic
validity does not require $\beta\ge0$, although the
lower-bound proposition requires $\beta>0$.
\item Periodicity ensures the integration-by-parts boundary
terms vanish. Analogous identities hold on $\mathbb R^d$
for fields with adequate decay and regularity, and on bounded
domains under suitably compatible boundary conditions.
\item The regularity in Theorem~\ref{thm:main} is a convenient
sufficient assumption for classical differentiation; extending
the identity to weaker solutions requires the corresponding
chain-rule or approximation argument.
\item A Lyapunov structure for an \emph{assumed dissipative}
effective amplitude equation does not prove that this
equation, its damping coefficient, or its nonlinear saturation
is derived from a separate Lorentzian higher-derivative action.
\item The opposite sign $+b(\Delta+k_\star^2)^2A$ with $b>0$
would grow high-frequency modes and is incompatible with
the damping assertion in \eqref{eq:bandshell}.
\end{enumerate}

\section{Conclusion}

The E14 functional generates the E13 flow exactly:
\[
\boxed{\quad
\partial_t A=-\frac{\delta\E}{\delta\overline A},
\qquad
\frac{\dd}{\dd t}\E[A(t)]
=-2\norm{\partial_t A(t)}_{L^2}^2\le0.
\quad}
\]
Thus E14 is proved as a Lyapunov functional for the specified
periodic band-pass gradient-flow model, with its interpretation
and limitations explicitly stated.

\end{document}